\documentclass[11pt,a4paper]{article}
\usepackage[utf8]{inputenc}
\usepackage[T1]{fontenc}
\usepackage[margin=2.5cm]{geometry}
\usepackage{booktabs}
\usepackage{amsmath}
\usepackage{siunitx}
\usepackage{multirow}
\usepackage{float}
\usepackage[autostyle]{csquotes}
\usepackage[hidelinks]{hyperref}
\usepackage{graphicx}

\sisetup{detect-weight=true}

\title{\bfseries Marocco FL thesis --- week 05/10 update}
\author{}
\date{}

\begin{document}
\maketitle
\vspace{-3.2em}
\begin{center}
\emph{Client-Driven Local Work Assignment in Battery-Constrained Federated Learning}
\end{center}
\vspace{1em}

%======================================================================
\section*{Summary}

\begin{itemize}
  \item SAGE-smart is now \textbf{client-driven}: the server picks \emph{who}
        trains and proposes a range of work; each client decides \emph{how much}
        to train ($E$, $B$) with its own up-to-date state.
  \item The client--server protocol follows the professor's note: we combine
        \textbf{option 2} (server-side estimate) and \textbf{option 3} (the
        client accepts or refuses the proposal, refusals are replaced).
        Option 1 (broadcast) was tested and discarded: too expensive.
  \item A message is no longer free: every radio activation costs about
        \SI{12.6}{\joule} (LTE, Huang et al.\ 2012). This applies to all
        algorithms.
  \item Every algorithm (SAGE, ESCS-SD, SAGE-smart) now runs with three battery
        models: linear, Peukert, datasheet (nm). Results: 33 runs, one seed,
        three $\beta$.
\end{itemize}

%======================================================================
\section{The problem: the server selects with old information}

The server chooses clients by their state of charge (SoC). But the SoC it knows
is the one the client reported the \emph{last time it took part}, possibly 10
rounds ago. In the meantime the client has consumed energy in idle, or has
been recharged, or its workload has changed. The server is therefore choosing
with stale information.

The professor's note proposes three ways to fix this:

\begin{center}
\begin{tabular}{lll}
\toprule
option & what happens & cost \\
\midrule
1 -- broadcast & the server asks \emph{every} client for its state & one message per client per round \\
2 -- estimate  & the server extrapolates the SoC from the last report & zero \\
3 -- accept / refuse & the selected client checks its own state and can refuse & one message per candidate \\
\bottomrule
\end{tabular}
\end{center}

%======================================================================
\section{How much does a message cost?}

With an energy model in J/MB only, a 1\,KB status message would cost about
\SI{0.002}{\joule}: asking the state would look free. On a phone it is not.
On LTE, every isolated transmission wakes the radio up: it pays a
\emph{promotion} to the connected state, and then the radio stays on for a
\emph{tail} before going back to idle. Values from Huang et al.\ (MobiSys 2012,
Table 3):

\[
  E_{\text{session}} =
  \underbrace{\SI{1.21}{\watt}\times\SI{0.26}{\second}}_{\text{promotion}}
  + \underbrace{\SI{1.06}{\watt}\times\SI{11.58}{\second}}_{\text{tail}}
  \approx \SI{12.6}{\joule}
\]

regardless of the message size. Consequences:

\begin{itemize}
  \item \textbf{Model exchange} (all algorithms): download and upload are two
        sessions, separated by the training. Cost per participant per round:
        \SI{4}{\joule} of payload $+ 2\times\SI{12.6}{\joule} \approx \SI{29}{\joule}$
        (it was \SI{4}{\joule}).
  \item \textbf{Control message} (proposal and answer): \SI{12.6}{\joule} if it is
        isolated, about \SI{0.002}{\joule} if the model download follows right
        after, because the radio is already on.
  \item \textbf{Delay}: each exchange delays the round by promotion $+$ RTT
        $\approx \SI{0.26}{\second} + \SI{0.07}{\second} = \SI{0.33}{\second}$.
\end{itemize}

%======================================================================
\section{Testing the three options}

All options were implemented and compared on a logic simulator (real battery
physics, fake training), 30 clients, 6 selected per round, 250 rounds.
\enquote{State error} is the mean gap between the SoC the server believes and
the real one.

\begin{center}
\begin{tabular}{lccc}
\toprule
 & state error & extra energy (250 rounds) & wasted downloads \\
\midrule
stale (no protocol) & \SI{0.94}{\percent} & 0 & -- \\
1 -- broadcast & 0 & \SI{21}{\watt\hour} ($\approx\SI{12}{\percent}$ of total) & 0 \\
2 -- estimate & \SI{0.89}{\percent} & 0 & 366$^{*}$ \\
2 $+$ 3 & \SI{0.88}{\percent} & \SI{1.7}{\watt\hour}$^{*}$ & 0$^{*}$ \\
\bottomrule
\end{tabular}

\smallskip
{\footnotesize $^{*}$ stress test with the reserve raised to 0.62, to force refusals.}
\end{center}

\paragraph{Option 1} gives perfect information, but each round the $\approx 24$
clients asked and then not selected pay a full radio session for nothing.

\paragraph{Option 2} is free but blind to recharges and to workload changes
(the note says it too: the SoC can be estimated, the load cannot). Without
recharge its error drops from \SI{0.44}{\percent} to \SI{0.28}{\percent}: almost
all of its error comes from recharges it cannot see.

\paragraph{Option 3} lets a client refuse \emph{before} downloading the model.
In the stress test, the estimate alone produced 366 clients that downloaded the
model and then had to withdraw; with $2+3$ the 484 refusals all happened before
the download, were replaced, and nobody withdrew.

%======================================================================
\section{Our choice: option 2 $+$ option 3}

\begin{itemize}
  \item \textbf{Option 2 ranks everybody for free.} For most clients the
        estimate is good enough, because the battery changes slowly (one
        epoch costs about \SI{0.2}{\percent} of the battery).
  \item \textbf{Option 3 refreshes only the clients that matter}: the
        candidates. They answer with their current state and can refuse
        before the model is sent.
  \item \textbf{Cost}: a client that accepts pays $\approx\SI{0.002}{\joule}$, one
        that refuses pays \SI{12.6}{\joule}. At most 2 waves of proposals per
        round ($\approx\SI{0.33}{\second}$ each).
  \item \textbf{No starvation}: a refusal counts as a contact, so its
        staleness bonus restarts from zero and it is not asked again at the
        next round.
\end{itemize}

What it implies: real Flower messages (type \texttt{query}, with a new handler
on the client); if nobody accepts after 2 waves the run ends as \enquote{pool
exhausted}; a small gap remains between the answer and the training (the
workload may change), so a client can still withdraw at training time.

%======================================================================
\section{How SAGE-smart works now, step by step}

\paragraph{0. Subscription (once).} The server records the static data of each
client: tier, capacity, speed, dataset size, initial SoC. From these it computes
the \textbf{round deadline} (median time of 5 epochs at $B=128$, about
\SI{100}{\second} here) and the \textbf{batches to propose}: the energy-optimal
batch of every device rounds to 32, so the proposal is $B\in\{32,64\}$.

\paragraph{1. Estimate (option 2).} For every client, with no messages:
\[
  \mathrm{SoC}_{\text{est}} = \mathrm{SoC}_{\text{last}}
  - \frac{P_{\text{idle}}(w_{\text{last}})\,\Delta t}{V_{\text{nom}}\,C},
  \qquad P_{\text{idle}}(w) = \SI{0.1}{\watt} + w\,(\SI{3}{\watt}-\SI{0.1}{\watt}).
\]

\paragraph{2. Ranking.}
\[
  \text{score} = a\,\mathrm{SoC}_{\text{est}} + b\,D + c\,\min\!\left(1, \tfrac{\text{rounds since last selection}}{25}\right)
\]
with $D$ the class divergence, $c=0.3$, and $(a,b)$ tuned per $\beta$. The
number of selected clients grows from 2 at round 1 to 6 from round 5 on.

\paragraph{3. Proposal (option 3).} Each candidate receives, without the model:
$E\in[E_{\text{lo}},E_{\text{hi}}]$ with $E_{\text{lo}}=2$ and
$E_{\text{hi}}=\min(10,\ E_{\text{deadline}},\ E_{\text{energy}})$, the batches
$\{32,64\}$, the deadline and the reserve (\SI{20}{\percent}).

\paragraph{4. Accept or refuse.} The client reads its current SoC, workload and
charging state, and accepts if it can do at least 2 epochs and stay above
\SI{20}{\percent}. Refusals are replaced by the next ones in the ranking.

\paragraph{5. The client decides (the core of the thesis).} After receiving the
model, the client:
\begin{enumerate}
  \item measures how useful it is: the loss $L$ of the global model on 256 local samples,
        $u=\min(1, L/\ln 10)$ ($\ln 10 \approx 2.30$ is the loss of random guessing);
  \item measures how much battery it can spend: $h=(\mathrm{SoC}-0.2)/0.8$ ($h=1$ when charging);
  \item chooses $E^{*} = E_{\text{lo}} + \operatorname{round}\big((E_{\text{hi}}-E_{\text{lo}})\,u\,h\big)$;
  \item goes down from $E^{*}$ until the predicted SoC \emph{at the end of the round}
        stays above \SI{20}{\percent} and the time, computed with the
        \emph{current} workload, fits the deadline. Among $\{32,64\}$ it takes
        the batch that leaves more charge.
\end{enumerate}
Example measured: with the same deadline, a client with no workload chooses
$E=5$; with \SI{10}{\percent} workload $E=4$; with \SI{30}{\percent}
workload $E=3$. The server could not know this: this is the \emph{late binding}.

\paragraph{6. Answer and aggregation.} The client sends the update and declares
$E$, $B$, the steps done, the initial and \emph{final SoC}, the loss before and
after. The server checks that $(E,B)$ is inside the proposal and aggregates
with FedNova, using the declared steps $\tau_i$:
\[
  x \leftarrow x + \tau_{\text{eff}} \sum_i p_i \frac{\Delta_i}{\tau_i},
  \qquad \tau_{\text{eff}}=\sum_i p_i\tau_i .
\]

\paragraph{7. Update.} The declared final SoC becomes the new starting point
for the estimate of step 1.

%======================================================================
\section{Three battery models for every algorithm}

The real battery of the simulator is always the datasheet model. What changes
is the model the \emph{algorithm} uses to estimate the SoC:

\begin{center}
\begin{tabular}{llll}
\toprule
 & linear (lin) & Peukert (peuk) & datasheet (nm) \\
\midrule
SAGE       & \texttt{sage} (the paper)   & \texttt{sage\_peuk}       & \texttt{sage\_soc} \\
SAGE-smart & \texttt{sage\_smart\_lin}   & \texttt{sage\_smart\_peuk} & \texttt{sage\_smart} \\
ESCS-SD    & \texttt{escs\_sd\_lin}      & \texttt{escs\_sd\_peuk}    & \texttt{escs\_sd} \\
\bottomrule
\end{tabular}
\end{center}

\begin{itemize}
  \item \textbf{nm}: the algorithm sees the true SoC.
  \item \textbf{peuk}: a fuel gauge using Peukert ($n=1.15$), following
        training, communication, idle and recharge.
  \item \textbf{lin} for \textbf{SAGE-smart}: a linear fuel gauge
        ($P\,t/(V_{\text{nom}}C)$), also following idle and recharge.
  \item \textbf{lin} for \textbf{SAGE and ESCS}: the accounting of the papers,
        $\mathrm{SoC}_0 - \text{Wh of the rounds the client trained}/C$.
        It does \emph{not} count idle and recharge.
\end{itemize}

How wrong each model is, compared with the real battery (predicted drain /
true drain, \SI{3500}{\milli\ampere\hour} cell):

\begin{center}
\begin{tabular}{lcc}
\toprule
 & linear & Peukert \\
\midrule
idle (\SI{0.1}{\watt})      & 0.96 -- 1.09 & \textbf{0.59 -- 0.67} \\
training (2 -- \SI{3}{\watt}) & 0.95 -- 1.09 & 0.92 -- 1.11 \\
\bottomrule
\end{tabular}
\end{center}

The linear model is almost right at every current. Peukert with $n=1.15$
underestimates the drain at low current: in idle it sees about \SI{40}{\percent}
less consumption than the real one.

%======================================================================
\section{Results}

One seed (2026), $\beta\in\{0.1, 0.5, 1.0\}$, 250 rounds. Mean over the three
$\beta$; per $\beta$ the number of depleted devices is shown in the second table.

\begin{center}
\begin{tabular}{llcccc}
\toprule
algorithm & battery & accuracy & final SoC & depleted & energy [Wh] \\
\midrule
\multirow{3}{*}{SAGE}       & lin (paper) & 0.530 & 0.33 & \textbf{11.7} & 277 \\
                            & peuk        & 0.515 & 0.38 & 2.3  & 276 \\
                            & nm          & 0.513 & 0.39 & 2.7  & 270 \\
\midrule
\multirow{3}{*}{SAGE-smart} & lin         & 0.620 & 0.57 & 0    & 147 \\
                            & peuk        & 0.631 & 0.57 & 0    & 149 \\
                            & nm          & \textbf{0.633} & 0.56 & 0 & 148 \\
\midrule
\multirow{3}{*}{ESCS-SD}    & lin (paper) & 0.583 & 0.33 & \textbf{11.7} & 284 \\
                            & peuk        & 0.545 & 0.29 & 9.3  & 315 \\
                            & nm          & 0.568 & 0.32 & 7.3  & 320 \\
\midrule
FedAvg  & -- & 0.577 & 0.18 & 18 & 322 \\
FedProx & -- & 0.568 & 0.18 & 18 & 322 \\
\bottomrule
\end{tabular}
\end{center}

\begin{center}
\begin{tabular}{llccc}
\toprule
 & & \multicolumn{3}{c}{depleted devices} \\
algorithm & battery & $\beta=0.1$ & $\beta=0.5$ & $\beta=1.0$ \\
\midrule
\multirow{3}{*}{SAGE}    & lin (paper) & 12 & 12 & 11 \\
                         & peuk        & 3  & 1  & 3  \\
                         & nm          & 4  & 0  & 4  \\
\midrule
SAGE-smart               & all three   & 0  & 0  & 0  \\
\midrule
\multirow{3}{*}{ESCS-SD} & lin (paper) & 12 & 11 & 12 \\
                         & peuk        & 13 & 8  & 7  \\
                         & nm          & 9  & 6  & 7  \\
\bottomrule
\end{tabular}
\end{center}

\subsection{Reading the results}

\paragraph{The link is on depleted devices: the more faithful the battery
state the algorithm uses, the fewer devices it depletes.} Accuracy stays almost
flat.

\paragraph{SAGE.} 11.7 depleted devices with the paper accounting, 2--3 with
Peukert or the true SoC; the first death moves from round 118--138 to round
194--244 (or never). The reason is not linearity: the paper accounting
\emph{ignores idle}, which is about half of all the energy consumed and, by the
end of the run, \SIrange{30}{38}{\percent} of SoC per device that the selector
does not see. It believes the devices are fuller than they are, keeps choosing
them and depletes them. Peukert and nm behave the same: the Peukert error is a
few percent, too small to change who gets selected. The slightly higher
accuracy of the paper version comes with 3--4 times more depleted devices (and
it also uses different weights and rules).

\paragraph{ESCS-SD.} Same direction: 11.7 $\rightarrow$ 9.3 $\rightarrow$ 7.3
depleted devices, and the first death comes later (e.g.\ $\beta=0.5$: round
101 $\rightarrow$ 117 $\rightarrow$ 125). Accuracy is equal at $\beta=0.5$,
growing at $\beta=1.0$ (0.646 $\rightarrow$ 0.654 $\rightarrow$ 0.667), noisy at
$\beta=0.1$ with one seed. ESCS still depletes 7 devices with the true SoC: its
filter only looks at the next round, and its rounds are long (5 epochs, no
deadline), so everybody consumes a lot in idle. Note: the linear arm seems to
use less energy only because dead devices stop consuming.

\paragraph{SAGE-smart.} No depleted device with any battery model, final SoC
$\approx 0.57$, and about \textbf{half the energy} of the others (\SI{148}{\watt\hour}
vs 270--\SI{320}{\watt\hour}). Here the battery model barely matters: at
$\beta=0.5$ and $1.0$ the three variants are within 0.007 of accuracy. This is
expected, for three reasons:
\begin{enumerate}
  \item its fuel gauges count idle too, so the linear one is wrong only by its
        non-linearity: about $-0.5\%$ of SoC at the end of the run (Peukert
        $+3$ to $+4\%$), far from the \SI{20}{\percent} reserve;
  \item the deadline makes rounds short: total idle energy is \SI{91}{\watt\hour}
        against 165--\SI{173}{\watt\hour} of the other algorithms, so batteries stay high;
  \item only at $\beta=0.1$, where the loss is higher and clients train more
        (3.4 epochs on average vs 2.2), do batteries get closer to the reserve:
        there the order is lin $<$ peuk $<$ nm in accuracy (0.498, 0.520, 0.535).
\end{enumerate}
The negotiation stayed almost idle: 0--5 refusals per run, no withdrawals,
control energy below \SI{0.02}{\watt\hour}.

%======================================================================
\newpage
\section{Figures}

\begin{figure}[H]
    \centering
    \includegraphics[width=0.75\linewidth]{3_fig5_pareto_soc_ball.png}
    \caption{Final SoC vs accuracy, all $\beta$ pooled. Larger marker = more depleted devices.}
\end{figure}

\begin{figure}[H]
    \centering
    \includegraphics[width=\linewidth]{3_figB_pareto_soc_models_ball.png}
    \caption{The same plot, one panel per algorithm: linear (red), Peukert (orange), datasheet (green).}
\end{figure}

\begin{figure}[H]
    \centering
    \includegraphics[width=0.75\linewidth]{3_fig4_soc_ball.png}
    \caption{Final SoC and depleted devices per algorithm.}
\end{figure}

\begin{figure}[H]
    \centering
    \includegraphics[width=0.75\linewidth]{3_fig2_accuracy_rounds_ball.png}
    \caption{Accuracy per round.}
\end{figure}

%======================================================================
\newpage
\section{Limits and next steps}

\begin{itemize}
  \item \textbf{One seed}: the trend on depleted devices is clear
        (11.7 $\rightarrow$ 2.5 for SAGE), the differences in accuracy are not.
        Run seeds 7 and 13.
  \item \textbf{Separate idle from linearity}: SAGE and ESCS linear arms ignore
        idle, so they do not isolate the battery model. Add a full linear fuel
        gauge for them too, as SAGE-smart already has.
  \item \textbf{Where does the SAGE-smart gain come from?} It changes several
        things at once (deadline, fewer epochs, $B=32$, FedNova, staleness,
        client decision). Ablations: without late binding, without deadline,
        without the loss term.
  \item \textbf{Fair amount of work}: baselines always do $E=5$, $B=128$;
        SAGE-smart about 2--3 epochs. Add a baseline with the same number of
        epochs and compare energy per accuracy.
  \item \textbf{A more dynamic world}: the workload changes once per round and
        the battery drains slowly, so the information only the client has is
        worth little and the negotiation is rarely triggered.
\end{itemize}

\end{document}
